\section{Background}

% --- 2.1 Accountability in Multi-Agent AI Systems ---

\subsection{Accountability in Multi-Agent AI Systems}

The problem of accountability in autonomous multi-agent systems has grown acute with the transition from single-agent pipelines to dynamic, recursive agent-spawning architectures. Early multi-agent systems literature treated accountability as a property of the individual agent --- an agent is accountable if its designers can explain its behavior \cite{wooldridge1995,stone2000}. That framing is insufficient when agents spawn sub-agents, delegate tasks across trust boundaries, and operate in environments where no single developer retains full situational awareness. As the delegation chain branches, loops, or creates cycles through recursive spawning, the provenance of any given action becomes a graph rather than a traceable chain. Mutable delegation chains permit consequential actions to proceed without any checkpoint where a human principal can intervene or verify.

Dijkstra's structured programming paradigm established a foundational principle relevant to this problem: systems should have statically verifiable structures that prevent silent entry into unintended states \cite{dijkstra1982}. In the control-flow context, unrestricted jumps permit the program to enter arbitrary states from which recovery is not guaranteed. The structured programming response --- making the set of reachable states statically enumerable --- was a fail-safe measure. Applied to agentic systems, the principle implies that mutable delegation chains constitute a structural accountability failure analogous to unrestricted jumps in program control flow. The BX3 Framework addresses this through its Fact Layer, which maintains an immutable authorization ledger independent of any agent process \cite{bx3framework2026}. The ledger records the human principal's identity, the scope of authorization, and the termination conditions for every spawn event, making authorization state tamper-evident across agent restarts.

A second dimension of the accountability problem is temporal. An agent may spawn a child that executes over hours or days. During that interval, the parent agent may be suspended, restarted, or replaced. The question of which human principal authorized a long-running task is not answered by a delegation certificate alone --- it requires a durable, immutable record that persists regardless of agent lifecycle events. The BX3 Framework's Forensic Ledger satisfies this requirement by anchoring every spawn event to a designated human anchor $h(n)$ via an immutable parent pointer $\alpha(n)$, ensuring that the escalation path $C(n)$ from any node $n$ to its human principal is traceable even after agent restarts \cite{bx3framework2026}.

% --- 2.2 Human-in-the-Loop Design ---

\subsection{Human-in-the-Loop Design}

The human-in-the-loop (HITL) paradigm is the dominant approach to maintaining human oversight of autonomous systems. Its canonical formulation traces to Wiener's \emph{Cybernetics}, which established that any system capable of purposeful action must include a feedback loop that reaches a human operator capable of modifying the system's goals \cite{wiener1948}. Wiener's formulation was prescriptive: systems without human oversight were not merely inefficient but ethically inadmissible, because they divorced decision-making from accountability. The human-in-the-loop was not an optional feature --- it was a constitutive requirement for any system that issues consequential directives.

The practical implementation of HITL in modern AI systems has diverged significantly from Wiener's original intent. Traditional HITL was a design-time construct: a human defined the system's goals, the system executed, and the human intervened only at clearly defined breakpoints. Contemporary LLM-based agents operate in continuous time with no natural breakpoints --- and critically, the decision of \emph{when} to escalate is itself made by the agent, not by the architecture. This creates a self-referential problem: if the agent decides when to call a human, the agent can defer escalation precisely when it is most needed. The HITL checkpoint becomes a procedural formality rather than a genuine oversight mechanism.

Bansal et al.\ provide empirical evidence for this failure mode in the context of human-AI collaboration \cite{bansal2019}. Their central finding is that HITL mechanisms improve team performance only when the human operator possesses an accurate mental model of the AI's decision-making process. When the AI's reasoning is opaque, or when the escalation interface lacks diagnostic context, humans default to either over-trusting or under-trusting the system --- both of which degrade safety outcomes. The content and format of information presented at the HITL checkpoint is not an ergonomic afterthought; it determines whether the checkpoint functions as a genuine corrective mechanism or as a ritual that provides psychological comfort without actual safety leverage \cite{bansal2019}.

This finding has direct consequences for the design of escalation mechanisms in agentic systems. A HITL interface that surfaces only a binary exception flag --- \emph{error occurred, human required} --- provides no diagnostic context for the human to evaluate the severity or nature of the failure. The BX3 Framework's Bailout Protocol addresses this by requiring that every escalation carry a complete diagnostic record: the exception classification, the operating range definition, the Bounds Engine's reasoning trace, and the resolution history of prior occurrences at that node \cite{bailoutprotocol2026}. This is a structural requirement for the HITL checkpoint to function as Wiener intended --- not merely as a flag, but as a meaningful feedback loop with genuine corrective power.

% --- 2.3 Fail-Safe Design in Safety-Critical Systems ---

\subsection{Fail-Safe Design in Safety-Critical Systems}

The concept of fail-safe design originates in control theory and aerospace engineering, where the cost of failure is measured in human life. A fail-safe system is one that transitions to a safe state upon the occurrence of any fault that is either unanticipated or outside the system's operational envelope. The canonical fail-safe property is \emph{graceful degradation}: a system that cannot complete its intended function must fail in a way that causes no harm and that preserves the maximum possible measure of system integrity.

Dijkstra's structured programming paradigm established a control-flow discipline directly applicable to fail-safe system design: programs should have statically verifiable structures that prevent the system from entering states that are neither designed nor intended \cite{dijkstra1982}. The \emph{go-to controversy} was, at its core, a fail-safe argument. Unrestricted jumps permit the program to enter arbitrary states from which recovery is not guaranteed. Structured programming proposed a discipline that made the set of reachable states \emph{a priori} enumerable --- and therefore inspectable and, crucially, safe to reason about. The analogy to agentic systems is direct: an unrestricted agent reasoning trace can enter arbitrary decision states from which recovery is not guaranteed, and the only architectural remedy is to bound the reachable state space.

Tristr's bounded rationality framework extends this argument to the design of complex organizational and computational systems \cite{tristr81}. The core thesis is that any system that attempts to operate beyond its defined bounds --- whether those bounds are computational, informational, or organizational --- will experience a form of systemic failure that is not recoverable by the system itself. The implication for agentic AI is direct: an agent that exceeds its provisioned operating range cannot be relied upon to recover itself, because the recovery action requires the very cognitive resources that are no longer available within the defined bounds. The only recoverable failure mode is one where the excess is detected and escalated before the agent's reasoning capacity is compromised \cite{tristr81}.

This logic maps directly onto the BX3 Framework's architectural commitment to bounded autonomy. The Purpose Layer defines the goal; the Bounds Layer enforces the operating envelope; the Fact Layer records and verifies every state transition. An agent that attempts to operate outside its provisioned bounds enters the Bailout Protocol: the system transitions to a safe state (human escalation) and stalls all consequential action until a human principal evaluates the exception and issues a directive. The Fact Layer's pre-commit hook enforces this atomically --- no machine actor, including the agent whose behavior triggered the escalation, can interpose or suppress the transition to the \texttt{ESCALATING} state \cite{bx3framework2026}. The architectural commitment is not to prevent agents from erring --- it is to ensure that when they do, the error cannot propagate silently.

% --- 2.4 The BX3 Framework: Three-Layer Model ---

\subsection{The BX3 Framework: Three-Layer Model}

The BX3 Framework is a structural architecture for accountable autonomous systems that synthesizes the accountability, HITL, and fail-safe requirements above into a concrete three-layer model. It decomposes every agent node into three immutable functional layers, each with a distinct and non-overlapping role in the accountability chain:

\begin{enumerate}
    \item \textbf{Purpose Layer.} The Purpose Layer encodes the intended goal state and the terminal success criteria for the agent's current task. It is the sole source of directives to the agent's action space. The Purpose Layer is human-authored --- either directly or through an authenticated delegation chain that traces to a human principal. No machine actor can modify the Purpose Layer's encoded intent \cite{bx3framework2026}.
    \item \textbf{Bounds Layer.} The Bounds Layer defines the operating envelope within which the agent is authorized to act. This includes resource constraints (time, memory, compute budget), scope constraints (which actions are permitted in the current context), and norm constraints (which categories of outcome are classified as unacceptable regardless of goal attainment). The Bounds Layer monitors the agent's reasoning trace against these constraints and triggers the Bailout Protocol when a constraint violation is detected or when the Bounds Engine's timeout $T_{\text{bounds}}$ expires without resolution \cite{bx3framework2026}.
    \item \textbf{Fact Layer.} The Fact Layer maintains the immutable authorization ledger and the forensic record of all state transitions at the node. It is the only layer with write access to the authorization ledger. The Fact Layer enforces the Bailout Protocol's state machine transitions directly --- bypassing agent-level code --- via a pre-commit hook that rejects any directive that would violate the safety properties established in the system. Every ledger entry is SHA-256 hash-chained, making the record tamper-evident \cite{bx3framework2026}.
\end{enumerate}

The three layers communicate directionally: the Purpose Layer issues directives downward; the Bounds Layer monitors upward and outward, detecting constraint violations; the Fact Layer records laterally and upward, maintaining the immutable audit trail. The Bailout Protocol is the enforced communication path from the Bounds Layer to the human principal, traversing the immutable parent chain $\alpha(n)$ from the originating node $n$ to its designated human anchor $h(n)$ \cite{bailoutprotocol2026}.

The significance of this architecture is that it separates concerns that are conflated in conventional agentic systems. In a monolithic agent, goal reasoning, constraint monitoring, and authorization recording are performed by the same cognitive process --- which means the agent can, in principle, modify its own authorization state or suppress its own escalation. The BX3 Framework prevents this by making the Fact Layer's authorization ledger read-only from the agent's perspective and by enforcing state transitions through a pre-commit hook that cannot be intercepted by any layer below it. This is structurally analogous to hardware-enforced privilege separation in operating systems, and it is similarly foundational to the system's safety properties.

The BX3 Framework's three-layer model provides the architectural substrate for the Compulsory Human Escalation theorem: for any node $n$ in the spawn tree, and for any exception state $e$ at $n$, the Bailout Protocol's escalation path $C(n)$ always reaches the human anchor $h(n)$ in finite time \cite{bx3framework2026}. The theorem follows from the immutability of the parent pointer $\alpha(n)$ and the finiteness of the spawn tree depth, both of which are enforced by the Fact Layer's pre-commit hook. It establishes that the BX3 Framework's accountability guarantee is not a property of agent behavior --- it is a property of the architecture itself, and it holds regardless of the reasoning strategy employed by any individual agent node.
